\documentclass[11pt]{article}
\usepackage[margin=1in]{geometry}
\usepackage{amsmath,amssymb,amsthm}
\usepackage[colorlinks=true,linkcolor=blue,citecolor=blue,urlcolor=blue]{hyperref}

\newtheorem{theorem}{Theorem}
\newcommand{\R}{\mathbb{R}}
\newcommand{\Cb}{C_b}

\title{Literature Resolution of the Translation-Invariant Min--Max Question}
\author{Run \texttt{fa\_banach\_001}, lane 7}
\date{August 11, 2026}

\begin{document}
\maketitle

\paragraph{Status.}
The selected question from Guillen--Schwab's 2016 paper is answered affirmatively, and in the exact requested form, by their later paper on Euclidean space.  This packet records the match and does not claim a new proof.

\section*{The source question}
The 2016 paper proves min--max representations for nonlinear operators satisfying the global comparison property.  In its concluding list of questions, the authors ask whether, when the underlying manifold is \(\R^d\) and the nonlinear operator \(I\) is translation invariant, the linear operators in the min--max representation can all be chosen translation invariant \cite[Section 1.5]{GS2016}.

Here translation invariance means that, for translations \(\tau_xu(y)=u(x+y)\),
\[
 I(\tau_xu,y)=I(u,x+y).
\]
Equivalently, the whole operator is recovered from the scalar functional
\(F(u)=I(u,0)\) by \(I(u,x)=F(\tau_xu)\).

\section*{The later theorem}
The later paper states the following result as Theorem 1.10 \cite[Theorem 1.10]{GS2020}.

\begin{theorem}[Guillen--Schwab]
Suppose
\[
 I:C_b^2(\R^d)\longrightarrow C_b^0(\R^d)
\]
is Lipschitz, has the global comparison property, and is translation invariant.  Then there are index sets and constants \(f_{ab}\), together with linear translation-invariant constant-coefficient L\'evy operators \(L_{ab}\), such that
\[
 I(u,x)=\min_a\max_b\bigl\{f_{ab}+L_{ab}(u,x)\bigr\}.
\]
Every linear operator appearing in the representation therefore commutes with translations.
\end{theorem}

This is precisely the affirmative conclusion requested in the selected source question.  The theorem imposes the natural Euclidean function-space and Lipschitz hypotheses under which the original min--max theory is formulated.

\section*{Why the proof resolves the structural issue}
The later proof does more than infer translation invariance indirectly.  It starts with the scalar functional
\[
 F(u)=I(u,0),
\]
applies the scalar min--max representation to \(F\), and identifies its supporting linear functionals with constant-coefficient L\'evy operators.  Translation invariance then supplies
\[
 I(u,x)=F(\tau_xu),
\]
so the same constant-coefficient operators give the representation at every point \(x\).  Thus the representing family is translation invariant by construction, exactly addressing the possible loss of symmetry raised in 2016 \cite[proof of Theorem 1.10]{GS2020}.

\section*{Scope of this resolution}
The 2016 article ends with several broader questions concerning manifolds, regularity, and alternative domains.  The literature theorem above resolves the selected Euclidean translation-invariance question only.  No claim is made here about the other questions in that list.

\section*{Verification record}
The source question was checked in the authors' 2016 arXiv version.  The statement and proof mechanism above were checked in arXiv:1812.09642v2, where the result is numbered Theorem 1.10.  The later paper was published in \emph{Nonlinear Analysis} 193 (2020), article 111468, DOI \href{https://doi.org/10.1016/j.na.2019.02.021}{10.1016/j.na.2019.02.021}.

\begin{thebibliography}{9}
\bibitem{GS2016}
N. Guillen and R. W. Schwab,
\emph{Min--max formulas for nonlocal elliptic operators},
arXiv:1606.08417 (2016),
\url{https://arxiv.org/abs/1606.08417}.

\bibitem{GS2020}
N. Guillen and R. W. Schwab,
\emph{Min--max formulas for nonlocal elliptic operators on Euclidean space},
Nonlinear Analysis 193 (2020), 111468;
arXiv:1812.09642v2,
\url{https://arxiv.org/abs/1812.09642}.
\end{thebibliography}

\end{document}
